% Dot batch 77; source SHA-256: 121eb37740a8f4c914610412e27e25ae7feec06b0d70a04fed3cbe3d4ee4db60
% Archived source: dot_updates_20261124/batches/batch077/EGA第五卷_除子族的固定点与转置.tex
% Original body lines: 24--128; integration map: integration/late_integration.json
\begingroup
\def\newscope{V-L}
\def\Div{\operatorname{Div}}
\def\pr{\operatorname{pr}}
\def\Supp{\operatorname{Supp}}
\def\Ass{\operatorname{Ass}}
\def\Spec{\operatorname{Spec}}
\def\OO{\mathcal O}
\egasection{3\quad 除子的线性系与到射影丛的态射}
设 \(D\) 是由 \(T\) 参数化的 \(X/S\) 上的除子族。以下均指有效除子。
记 \(\pr_1:X\times_S T\to X\) 为第一投影。
若 \(\pr_1^{-1}(x)\subseteq \Supp D\) 在集合意义下成立，就称 \(x\in X\)
为该除子族的一个\emph{固定点}。因此，若用 \(Z\) 表示固定点集，则
\[
 X\setminus Z
  =\pr_1\bigl((X\times_S T)\setminus\Supp D\bigr).
\]
所以，若 \(T\to S\) 普遍开（例如平坦且局部有限表示），则 \(Z\) 闭。
若 \(Z=\varnothing\)，称此除子族\emph{无固定点}。
若 \(Z\) 闭，则 \(X\setminus Z\) 是 \(X\) 中最大的开集，使给定的除子族
限制到此开集后、仍以 \(T\) 参数化的诱导除子族没有固定点。\cmt{L18-077-01}{法稿手写
\emph{ensemblistement} 明确指支集的包含，而不要求闭子概形的包含。
\(\pr_1\) 是 \(T\to S\) 的基变换，所以普遍开用于非固定点集的开性。
法文物理126页先前划销的“除子族”草稿，另见本版“第三篇删划草稿补译”中对应页段，作为非定稿材料保存；仍未辨的笔迹在那里单列。}

若 \(D\) 无固定点，且 \(T\to S\) 平坦、局部有限表示，纤维满足 \((S_1)\)
并且几何不可约，则相对于第一投影 \(\pr_1\)，\(D\) 也是 \(X\) 上的相对除子。
确实，在 \(z\in\Supp D\) 附近，\(D\) 由一个方程 \(\phi=0\) 定义。
令 \(x=\pr_1(z)\)，并令 \(s\) 为 \(x\) 在 \(S\) 中的像。
在纤维 \(T_s\times_{k(s)}k(x)\) 上，该方程诱导的局部方程 \(\phi_x\)
在 \(z\) 处不是幂零元。否则，\(\pr_1^{-1}(x)\cap\Supp D\)
会包含这条纤维上 \(z\) 的一个开邻域，从而因纤维不可约而包含整个纤维；
这就把 \(x\) 变成了固定点，与假设矛盾。
又因 \(T_s\times_{k(s)}k(x)\) 不可约并满足 \((S_1)\)，
\(\phi_x\) 在 \(z\) 处是 \(\OO_{T_s\times_{k(s)}k(x),z}\) 的非零因子。
这样得到由 \(X\) 参数化的 \(T/S\) 上的一族除子，也就是一个态射\cmt{L18-077-02}{在局部诺特的 \((S_1)\) 概形中，关联点都是不可约分支的泛点；
不可约时，零因子只来自唯一的极小素理想，即幂零根。
故这里不需要把几何不可约加强成几何整。
有限型域代数的 \((S_1)\) 性质经过此处的域扩张保持；可用有限生成域扩张
给出的平坦映射有 Cohen--Macaulay 纤维，再用
\href{https://stacks.math.columbia.edu/tag/0336}{平坦上升判据}。
从纤维上的正则性回到相对 Cartier 除子，见
\href{https://stacks.math.columbia.edu/tag/056P}{相对有效 Cartier 除子一节}。}
\[
 X\longrightarrow\Div_{T/S}.
\]

一般而言，\(X/S\) 上的除子族可能有固定点。在上述对 \(T/S\) 纤维的假设下，
把 \(X\) 换成 \(X\setminus Z\)，即可得到由 \(X\setminus Z\) 参数化的
\(T/S\) 上的除子族，也就是
\[
 X\setminus Z\longrightarrow\Div_{T/S}.
\]
而且上述证明说明，\(X\setminus Z\) 正是使 \(D|_{U\times_S T}\)
成为相对于 \(U\) 的除子的最大开集 \(U\subseteq X\)；
换言之，其转置 \({}^{t}D\) 限制后是一族由 \(U/S\) 参数化的 \(T/S\) 上的除子。\cmt{L18-077-03}{法英两稿在“把 \(X\) 换成……”处均印作 \(X-T\)，
但本段前后的箭头与最大开集都明载 \(X-Z\)，且英译者也指出此误。本文订正。
这里最大性的逆向使用非空纤维：若整条非空纤维落在支集里，
在其泛点局部方程幂零，不能是非零因子；若允许空纤维，普通固定点定义还会因空包含而成立，而相对除子条件在空纤维上自动成立。}

\medskip
也可顺便指出：若 \(X\) 和 \(T\) 均在 \(S\) 上平坦、局部有限表示，
且纤维满足 \((S_1)\) 并且几何不可约，则转置 \(D\mapsto{}^{t}D\)
给出如下两类对象的双射：一类是由 \(T\) 参数化的 \(X/S\) 上无固定点的除子族；
另一类是由 \(X\) 参数化的 \(T/S\) 上无固定点的除子族。\cmt{L18-077-04}{英译把双射右侧的
\(T/S\) 错作 \(T/X\)；法稿如正文。两边的族都位于同一个纤维积上，只交换被参数化的一方和参数一方。}

若想在此陈述中去掉对 \(X/S\) 和 \(T/S\) 的纤维所作的特殊假设，
宜把“固定点”替换成“\emph{广义固定点}”。
所谓 \(D\) 的广义固定点，是使 \(D\) 不在 \(\pr_1^{-1}(x)\)
的所有点上都为 \(X\) 上相对除子的点 \(x\in X\)；即至少有一个点不满足相对除子条件。
设 \(W\subseteq X\times_S T\) 是 \(D\) 为相对于 \(X\) 的除子的开域。
则广义固定点集是
\[
 Z'=\pr_1\bigl((X\times_S T)\setminus W\bigr).
\]
若 \(T\to S\) 固有，\(Z'\) 因此为 \(X\) 的闭子集。
不论是否固有，都可在 \(X\setminus\overline{Z'}\) 上取得一族由该开集
参数化的 \(T/S\) 上的除子。\cmt{L18-077-05}{法稿此处确有 \(\overline{Z'}\) 上方的横线；
英稿不带闭包，而且将“若固有”写成类似已假定固有的表述。
固有性通过基变换使 \(\pr_1\) 闭，才可直接用 \(X\setminus Z'\) 作开底；
没有该保证时，最大开安全域是 \(X\setminus\overline{Z'}\)。}
\(T/S\) 的纤维满足 \((S_1)\) 且几何不可约这一假设，正是用来保证
\(Z=Z'\)，即普通固定点与广义固定点相同的。\cmt{L18-077-06}{仍按非空、平坦且局部有限表示
的纤维情形理解。去掉不可约性时，两个集合确可不同：
取 \(X=\mathbb A^1_k\)、\(T=\Spec k\amalg\Spec k\)，两参数处的除子
分别取原点与空除子。原点不是普通固定点，但在第一个参数分支上，
方程限制到原点纤维为零，故是广义固定点。}


从几何上说，为简便起见，设 \(S=\Spec k\)，\(k\) 是代数闭域。
对于平坦、局部有限表示的 \(T/S\)，可以通过取几何纤维来作这种化简。
说 \(x\in X(k)\) 是普通固定点，就是说每个 \(t\in T(k)\) 都满足
\(x\in\Supp D_t\)；说它是广义固定点，则是说存在 \(T\) 的一个关联素循环 \(T'\)，
使每个 \(t\in T'(k)\) 都满足 \(x\in\Supp D_t\)。\cmt{L18-077-07}{这里的关联素循环
以关联点的闭包为支集，可包括嵌入关联点所给的闭包。
在 \(T\) 局部有限型于代数闭域时，闭点足以检测其闭子集；
局部方程在某个关联点消失，表示它在该处不是非零因子。}

\medskip
\noindent\textbf{补记一项遗漏。}
固定点集 \(Z\) 的形成与 \(S\) 上的基变换相容。
另一方面，若 \(X\setminus Z\) 为开集（例如 \(T/S\) 平坦且局部有限表示），
且 \(T\to S\) 满射，则 \(X\setminus Z\) 在 \(X\) 中相对于 \(S\) 普遍概形稠密。
这是第11节的结果和下列事实的推论：对每个 \(s\in S\)，
\(Z_s\) 不含 \(X_s\) 中任何与 \(\OO_{X_s}\) 关联的点，
因为 \(X_s\) 上有效 Cartier 除子的支集不含这样的点。\cmt{L18-077-08}{满射限定是法文物理129页的
手写补入，英稿没有。若 \(T_s\) 为空，整个 \(X_s\) 都因空包含而成为普通固定点，
所以不能省略这一限定。对任意 \(s\)，选取 \(t\in T_s\)，再扩域到 \(k(t)\)：
有效 Cartier 除子不能包含任何关联点，且每个旧关联点在此域扩张后都有新关联点在其上方。
另保留应用 \href{https://numdam.org/item/PMIHES_1966__28__5_0/}{EGA IV 第11节}
纤维判据所需的环境假设，例如 \(X/S\) 平坦、局部有限表示。
这里只谈 \(Z\) 作为点集的基变换，不凭空赋予其一个闭子概形结构。}
\par\endgroup
